% ============================================================
%  documentation.tex  —  bookcover.cls documentation
%  Compile with LuaLaTeX (or XeLaTeX).
% ============================================================

\documentclass[12pt, a4paper]{article}

% ── Packages ─────────────────────────────────────────────────
\usepackage{fontspec}
\usepackage{geometry}
\usepackage{hyperref}
\usepackage{listings}
\usepackage{xcolor}
\usepackage{booktabs}
\usepackage{parskip}
\usepackage{titlesec}
\usepackage{tcolorbox}
\usepackage{enumitem}
\usepackage{microtype}

\geometry{margin=2cm}

% ── Fonts ────────────────────────────────────────────────────
\setmainfont{Latin Modern Roman}
\setmonofont{Latin Modern Mono}

% ── Colours ──────────────────────────────────────────────────
\definecolor{bcBlue}    {HTML}{1B6B7B}
\definecolor{bcGold}    {HTML}{D4A96A}
\definecolor{bcGray}    {HTML}{F5F5F5}
\definecolor{bcDarkGray}{HTML}{444444}
\definecolor{bcGreen}   {HTML}{2E7D32}
\definecolor{bcPurple}  {HTML}{6A1B9A}

% ── Code block style ─────────────────────────────────────────
\lstdefinestyle{latex}{
  language          = [LaTeX]TeX,
  backgroundcolor   = \color{bcGray},
  basicstyle        = \ttfamily\small,
  keywordstyle      = \color{bcBlue}\bfseries,
  commentstyle      = \color{bcGreen}\itshape,
  stringstyle       = \color{bcPurple},
  breaklines        = true,
  breakatwhitespace = true,
  frame             = single,
  framerule         = 0.4pt,
  rulecolor         = \color{bcGold},
  xleftmargin       = 6pt,
  xrightmargin      = 6pt,
  aboveskip         = 8pt,
  belowskip         = 8pt,
  tabsize           = 2,
  showstringspaces  = false,
}
\lstset{style=latex}

% ── Section headings ─────────────────────────────────────────
\titleformat{\section}
  {\Large\bfseries\color{bcBlue}}
  {\thesection}{1em}{}[\vspace{-0.5cm}\color{bcGold}\rule{\linewidth}{0.8pt}]

\titleformat{\subsection}
  {\large\bfseries\color{bcDarkGray}}
  {\thesubsection}{1em}{}

\titleformat{\subsubsection}
  {\normalsize\bfseries\color{bcDarkGray}}
  {\thesubsubsection}{1em}{}

% ── Hyperlinks ───────────────────────────────────────────────
\hypersetup{
  colorlinks = true,
  linkcolor  = bcBlue,
  urlcolor   = bcBlue,
  citecolor  = bcBlue,
  pdftitle   = {bookcover.cls — Documentation},
  pdfauthor  = {bookcover},
}

% ── Warning box ──────────────────────────────────────────────
\newtcolorbox{warning}{
  colback  = bcGold!15,
  colframe = bcGold,
  boxrule  = 0.6pt,
  arc      = 3pt,
  left = 8pt, right  = 8pt,
  top  = 6pt, bottom = 6pt,
  fontupper = \small,
}

% ── Note box ─────────────────────────────────────────────────
\newtcolorbox{note}{
  colback  = bcBlue!8,
  colframe = bcBlue,
  boxrule  = 0.6pt,
  arc      = 3pt,
  left = 8pt, right  = 8pt,
  top  = 6pt, bottom = 6pt,
  fontupper = \small,
}

% ── Inline helpers ───────────────────────────────────────────
\newcommand{\cmd}[1]{\texttt{\textbackslash #1}}
\newcommand{\env}[1]{\texttt{#1}}
\newcommand{\opt}[1]{\texttt{#1}}

% ════════════════════════════════════════════════════════════
\begin{document}
% ════════════════════════════════════════════════════════════

% ── Cover page ───────────────────────────────────────────────
\begin{titlepage}
  \centering
  \vspace*{3cm}
  {\fontsize{36}{40}\selectfont\bfseries\color{bcBlue}
    bookcover\texttt{.cls}\par}
  \vspace{0.5cm}
  {\color{bcGold}\rule{0.5\linewidth}{1.2pt}\par}
  \vspace{0.8cm}
  {\Large Documentation — User Manual\par}
  \vspace{0.4cm}
  {\normalsize \LaTeX{} class for generating print-ready book covers\par}
  \vfill
  {\small Compile with \textbf{Lua\LaTeX} or \textbf{Xe\LaTeX}\par}
  \vspace{1cm}
  {\small v2.0 — 2026\par}
\end{titlepage}

% ── Table of contents ────────────────────────────────────────
\tableofcontents
\newpage

% ════════════════════════════════════════════════════════════
\section{Installation}
% ════════════════════════════════════════════════════════════

Place \texttt{bookcover.cls} in the same folder as your \texttt{.tex} file,
or in any directory on your \LaTeX{} search path.
On Overleaf, simply upload the file to your project.

% ════════════════════════════════════════════════════════════
\section{Class options}
% ════════════════════════════════════════════════════════════

\begin{lstlisting}
\documentclass[
  width       = 148,    % book page width in mm        (default: 148)
  height      = 210,    % book page height in mm       (default: 210)
  margin      = 10,     % content margin in mm; groove side gets
                        % margin + groove automatically (default: 10)
  bleed       = 20,     % bleed extension in mm        (default: 20)
  groove      = 8,      % binding hinge width in mm    (default: 8)
  paperwidth  = 420,    % output sheet width in mm     (default: 420)
  paperheight = 297,    % output sheet height in mm    (default: 297)
  pages       = 300,    % number of pages (sets spine width)
  grammage    = 75,     % paper weight in g/m²         (default: 75)
  colorFront  = 1B6B7B, % front panel colour, hex without #
  colorBack   = 1B6B7B, % back panel colour,  hex without #
  colorSpine  = D4A96A, % spine colour,       hex without #
]{bookcover}
\end{lstlisting}

\subsection{Supported grammage values}

\begin{center}
\begin{tabular}{cc}
  \toprule
  \textbf{Grammage (g/m²)} & \textbf{Thickness per sheet} \\
  \midrule
  56  & 0.07 mm \\
  60  & 0.08 mm \\
  75  & 0.10 mm \\
  90  & 0.12 mm \\
  115 & 0.15 mm \\
  \bottomrule
\end{tabular}
\end{center}

% ════════════════════════════════════════════════════════════
\section{Sheet anatomy}
% ════════════════════════════════════════════════════════════

The full artwork is centred on the output sheet. Reading left to right:

\begin{center}
\ttfamily
|{\normalfont←}bleed{\normalfont→}|{\normalfont←}\ back\ {\normalfont→}|{\normalfont←}groove{\normalfont→}|{\normalfont←}spine{\normalfont→}|{\normalfont←}groove{\normalfont→}|{\normalfont←}\ front\ {\normalfont→}|{\normalfont←}bleed{\normalfont→}|
\end{center}

\begin{description}[style=nextline, leftmargin=2cm, labelwidth=1.8cm]
  \item[\opt{bleed}]  Colour extends beyond the trim edge so no white border appears after guillotining.
  \item[\opt{back}]   Back cover panel (\texttt{width} mm wide).
  \item[\opt{groove}] Hinge recess for hardcover binding. Filled with the panel colour; text is kept clear automatically
  \item[\opt{spine}]  Width calculated from \opt{pages} and \opt{grammage}.
  \item[\opt{front}]  Front cover panel (\texttt{width} mm wide).
\end{description}

% ════════════════════════════════════════════════════════════
\section{Spine width calculation}
% ════════════════════════════════════════════════════════════

\begin{center}
  \texttt{spine\_width = (pages × mm\_per\_sheet) / 2}
\end{center}

The division by 2 accounts for two pages per physical sheet. Example:

\begin{center}
  \texttt{600 pages × 0.10 mm/sheet ÷ 2 = 30 mm}
\end{center}

% ════════════════════════════════════════════════════════════
\section{Basic structure}
% ════════════════════════════════════════════════════════════

\begin{lstlisting}
\documentclass[
  pages=600, 
  grammage=75,
  colorFront=1B6B7B, 
  colorSpine=D4A96A]
{bookcover}

\begin{front}
  % front cover content
\end{front}

\begin{back}
  % back cover content
\end{back}

\begin{spine}
  % spine content
\end{spine}

\begin{document}
  \makecover
\end{document}
\end{lstlisting}

\begin{note}
\cmd{makecover} is the only command needed inside\env{document}.
\end{note}

% ════════════════════════════════════════════════════════════
\section{Environments}
% ════════════════════════════════════════════════════════════

\subsection{front}

Defines the content of the front (right-hand) cover panel.

\begin{itemize}[noitemsep]
  \item Live area width: \texttt{width} mm (groove is included in the
        underlying minipage but hidden from text by the automatic margin).
  \item Automatic margins: \opt{margin}+\opt{groove} on the left (hinge side),
        \opt{margin} on the right.
  \item Use \cmd{colorblock} for full-width colour bands that span the groove.
\end{itemize}

\begin{lstlisting}
\begin{front}
  \centering\color{white}
  {\Huge\bfseries BOOK TITLE}
  \vfill
  {\large\bfseries AUTHOR NAME}
  \vspace{4mm}
  \colorblock{D4A96A}{10mm}{\centering\color{1B6B7B} Publisher}
\end{front}
\end{lstlisting}

\subsection{back}

Defines the content of the back (left-hand) cover panel.

\begin{itemize}[noitemsep]
  \item Live area width: \texttt{width} mm.
  \item Automatic margins: \opt{margin} on the left,
        \opt{margin}+\opt{groove} on the right (hinge side).
\end{itemize}

\begin{lstlisting}
\begin{back}
  \vspace{10mm}
  \color{white}
  \setlength{\parskip}{14pt}
  {\normalsize Book description goes here.}
  \vfill
  {\centering\includegraphics[width=0.25\linewidth]{logo.png}\par}
  \hfill{\color{gold}\small\bfseries ISBN 978-0-000-00000-0}
  \vspace{1cm}
\end{back}
\end{lstlisting}

\subsection{spine}

Defines the content of the spine strip.

\begin{itemize}[noitemsep]
  \item Width is calculated automatically from \opt{pages} and \opt{grammage}.
  \item No automatic margins as the spine may be too thin — layout is fully manual.
  \item Use \cmd{spinerotated} to rotate text top-to-bottom.
  \item Use \cmd{centering} and \cmd{vfill} to distribute elements vertically.
  \item \cmd{colorblock} is \textbf{not} available inside \env{spine}.
\end{itemize}

\begin{lstlisting}
\begin{spine}
  \centering
  \vspace{4mm}
  \color{teal}
  \spinerotated{\normalsize\bfseries Author Name}\par
  \vfill
  \spinerotated{\large\bfseries Book Title}\par
  \vfill
  \spinerotated{\small 4th ed.}\par
  \vfill
  {\normalsize\bfseries Vol.\\1\par}
  \vspace{4mm}
\end{spine}
\end{lstlisting}

\begin{warning}
\textbf{Spine width and readability:} if the spine is narrower than
approximately 12 mm, consider using only the title — or omitting text
entirely — to avoid overcrowding.
\end{warning}

% ════════════════════════════════════════════════════════════
\section{Commands}
% ════════════════════════════════════════════════════════════

\subsection{\cmd{colorblock}}

Draws a full-panel-width colour band inside \env{front} or \env{back}.
Unlike a plain \cmd{colorbox}, it deliberately escapes the automatic content
margins so the band always spans the complete panel width, including the
groove area.

\begin{lstlisting}
\colorblock{color}{height}{content}

% Parameters:
%   color   -- xcolor name or HTML hex without #
%   height  -- band height (any TeX length: mm, pt, \textheight...)
%   content -- any LaTeX content; receives full panel \linewidth

% Examples:
\colorblock{D4A96A}{0.25\textheight}{
  \centering\color{1B6B7B}
  {\Huge\bfseries\MakeUppercase{Book Title}}
}

\colorblock{D4A96A}{10mm}{
  \centering\color{1B6B7B}{\bfseries Publisher Name}
}
\end{lstlisting}

\begin{note}
Content inside \cmd{colorblock} is vertically centred. If you need top or bottom alignment, use a \cmd{minipage} with the desired \texttt{[t]} or \texttt{[b]} to wrap the content argument.
\end{note}

\subsection{\cmd{spinerotated}}

Rotates content $-90°$ (top-to-bottom, standard book spine direction) and centres it across the full spine width.

\begin{lstlisting}
\spinerotated{text}
\spinerotated{\large\bfseries Book Title}\par
\vfill
\spinerotated{\normalsize Author Name}\par
\end{lstlisting}

\begin{warning}
\textbf{Do not} wrap multiple elements in a single \cmd{spinerotated} call if you intend to use \cmd{vfill} between them — the fill will not work inside a rotated box. Call \cmd{spinerotated} once per element instead, and do not forget the \cm{par} at the end.
\end{warning}

\subsection{\cmd{makecover}}

Triggers the full TikZ render of the cover sheet. Must be the \textbf{only content} inside \env{document}.

\begin{lstlisting}
\begin{document}
  \makecover
\end{document}
\end{lstlisting}

\cmd{makecover} performs the following steps in order:

\begin{enumerate}[noitemsep]
  \item Converts all lengths to unitless mm values for \texttt{pgfmath}.
  \item Computes horizontal and vertical zone boundaries.
  \item Fills background colours (back, front, spine) with a 0.2 mm overlap between adjacent fills to prevent hairline seams in PDF viewers.
  \item Draws crop marks at the four corners of the trim box.
  \item Draws dimension annotations below and to the right of the artwork (for reference only — you may remove before sending to press).
  \item Places content nodes for back, spine, and front panels inside clipped scopes so overflow cannot bleed across zones.
\end{enumerate}

% ════════════════════════════════════════════════════════════
\section{Automatic margin behaviour}
% ════════════════════════════════════════════════════════════

The \opt{margin} class option sets a horizontal content margin that is
applied automatically to \env{front} and \env{back}. The groove-adjacent
side always receives \opt{margin}+\opt{groove} to keep text clear of the
binding hinge:

\begin{center}
\begin{tabular}{lll}
  \toprule
  \textbf{Panel} & \textbf{Left margin} & \textbf{Right margin} \\
  \midrule
  front & margin + groove & margin \\
  back  & margin          & margin + groove \\
  \bottomrule
\end{tabular}
\end{center}

\cmd{colorblock} bypasses these margins intentionally so that colour bands
span the full panel. All other content respects them.

% ════════════════════════════════════════════════════════════
\section{Dimension annotations}
% ════════════════════════════════════════════════════════════

\cmd{makecover} automatically draws labelled arrows below and to the right
of the artwork showing:

\begin{itemize}[noitemsep]
  \item bleed, back, spine, front, bleed widths (horizontal)
  \item groove widths (horizontal, labelled below the arrow line)
  \item cover height (vertical, on the right)
\end{itemize}

These annotations are \textbf{not} part of the print-ready output.
To remove them before sending to press, comment out the annotation block
in \cmd{bc@drawall} (step 3 in the source), or simply crop the PDF to the
trim box using your prepress software.

% ════════════════════════════════════════════════════════════
\section{Complete example}
% ════════════════════════════════════════════════════════════

\begin{lstlisting}
\documentclass[
  margin      = 10,
  pages       = 600,
  grammage    = 75, 
  colorFront  = 1B6B7B,
  colorBack   = 1B6B7B,
  colorSpine  = D4A96A
]{bookcover}

\usepackage[english]{babel}  % hyphenation patterns
\usepackage{graphicx}        % \includegraphics
\usepackage{ragged2e}        % \justifying
\usepackage{fontspec}        % requires XeLaTeX or LuaLaTeX.
    \setmainfont{TeX Gyre Adventor}
    \newfontfamily\TitleFont{TeX Gyre Adventor}[LetterSpace=10.0]
\usepackage{xcolor}
    \definecolor{teal} {HTML}{1B6B7B}  % matches colorFront / colorBack
    \definecolor{gold} {HTML}{D4A96A}  % matches colorSpine
    \definecolor{white}{rgb}{1,1,1}

%% ============================================================
%% TYPOGRAPHIC SCALE
%% ============================================================
\def\hugefont  {\fontsize{48}{56}}  % main cover title
\def\largefont {\fontsize{30}{36}}  % spine title
\def\mediumfont{\fontsize{20}{24}}  % subtitle, edition
\def\normalfont{\fontsize{16}{19}}  % authors, publisher, description
\def\smallfont {\fontsize{12}{14}}  % spine edition, ISBN

%% ============================================================
%% BOOK METADATA
%% ============================================================
%% --- Front cover ---
\def\booktitle{Discrete\\Mathematics}
\def\booksubtitle{Finally Proving P\,=\,NP (We Think)}
\def\bookedition{$\infty$th edition}
\def\coverimage{cover_graph.pdf}
\def\bookauthors{%
  Al G. Bra
  \quad\textbullet\quad
  Kay O. S. Theory%
}
\def\bookpublisher{$\mathcal{O}(1)$ Press}

%% --- Back cover ---
\def\bookdescription{%
  This textbook provides a rigorous introduction to discrete
  mathematics. All proofs are straightforward and left as an
  exercise to the reader.

  The second edition features twice the content of the first,
  achieved via the Banach--Tarski decomposition method.
  The publisher has been informed.

  The third part covers graph theory, including a definitive
  solution to the Bridges of K\"{o}nigsberg problem.
  (We added a bridge. You're welcome.)%
}
\def\publisherlogo{logo.png}
\def\bookisbn{ISBN 978-0-314-15926-5}

%% --- Spine ---
\def\spinetitle{Discrete Mathematics}
\def\spineauthors{Bra \textbullet\ Theory}
\def\spineedition{$\infty$th ed.}
\def\spinevolume{Vol.\\$n$}

%% ============================================================
%% FRONT COVER
%% ============================================================
\begin{front}

  %% Title and Subtitle — full-width band at the top
  \colorblock{gold}{0.25\textheight}{%
    \centering
    \color{teal}

    {\hugefont\TitleFont\bfseries\selectfont
      \MakeUppercase{\booktitle}\par}

    \vspace{2mm}
    {\rule{0.85\textwidth}{2pt}\par}
    \vspace{2mm}

    {\mediumfont\bfseries\selectfont
      \MakeUppercase{\booksubtitle}\par}
  }

  \centering
  \color{gold}

  %% Edition — pushed to the right margin with \hfill.
  \vspace{5mm} 
  \hfill{\mediumfont\itshape\selectfont\bookedition}

  %% Cover image — fills available vertical space between edition line and author names.
  \vfill 
  \includegraphics[width=\linewidth]{\coverimage}
  
  %% Author names — bold, uppercased, centred.
  \vfill 
  {\normalfont\bfseries\selectfont\MakeUppercase{\bookauthors}}

  %% Publisher strip — another full-width band at the bottom.
  \vspace{4mm} 
  \colorblock{gold}{10mm}{%
    \centering
    \color{teal}
    {\normalfont\bfseries\selectfont\bookpublisher}
  }

\end{front}

%% ============================================================
%% BACK COVER
%% ============================================================
\begin{back}

  \vspace{10mm}
  \color{white}
  \setlength{\parindent}{0pt}
  \setlength{\parskip}{14pt}     % breathing room between paragraphs

  %% Book description — justified text using ragged2e's \justifying,
  {\normalfont\selectfont\justifying\bookdescription}

  \vfill

  %% Publisher logo — centred, sized to 25% of the text width.
  {\centering\includegraphics[width=0.25\linewidth]{\publisherlogo}\par}

  %% ISBN — right-aligned, at the very bottom.
  \hfill{\color{gold}\smallfont\bfseries\selectfont\bookisbn}

  \vspace{1cm}
\end{back}

%% ============================================================
%% SPINE
%% ============================================================
\begin{spine}
  \centering
  \vspace{2mm}
  \color{teal}

  %% Authors at the top
  \spinerotated{{\normalfont\bfseries\selectfont\spineauthors}}\par

  %% Title in the centre — largest element on the spine
  \vfill 
  \spinerotated{{\largefont\bfseries\selectfont\spinetitle}}\par

  %% Edition below the title
  \vfill 
  \spinerotated{{\smallfont\bfseries\selectfont\spineedition}}\par

  %% Volume badge — upright text at the bottom of the spine
  \vfill 
  {\centering\normalfont\bfseries\selectfont\spinevolume\par}

  \vspace{2mm}

\end{spine}

%% ============================================================
%% DOCUMENT BODY
%% \makecover is the only content needed
%% ============================================================
\begin{document}
\makecover
\end{document}
\end{lstlisting}

\end{document}